\documentclass[10pt,journal,comsoc]{IEEEtran}
\usepackage{amsmath,amsfonts}
\usepackage{algorithm}
\usepackage{array}
\usepackage{amssymb}
\usepackage{amsthm}
\usepackage[caption=false,font=normalsize,labelfont=sf,textfont=sf]{subfig}
\usepackage{textcomp}
\usepackage{stfloats}
\usepackage{url}
\usepackage{verbatim}
\usepackage{float}
\usepackage{cite}
\usepackage{color}
\usepackage{multicol}
\usepackage{multirow}
\usepackage{graphicx}
\usepackage{mathtools}
\newtheorem{definition}{Definition}
\usepackage{makecell}
\usepackage{algpseudocode}
\usepackage{tabularx}
\usepackage{booktabs}
% Result figures (T1..T7) are written by tools/plot_thesis.py to
%   host_data/results/thesis_figures/ and mirrored to thesis/figures/.
\graphicspath{{./Figures/}{../thesis/figures/}{../host_data/results/thesis_figures/}}
\renewcommand{\arraystretch}{1.2}
\hyphenation{op-tical net-works semi-conduc-tor IEEE-Xplore}
\setlength{\emergencystretch}{1em}
\hbadness=10000
\vbadness=10000

\begin{document}

\title{Changed Decisions, Unchanged Delivery: When Do Adversarial Observation
Attacks Actually Harm MADDPG Routing?}

\author{
  Pedro Amaral,~\IEEEmembership{Member,~IEEE,}%
  \thanks{Pedro Amaral is with the Departamento de Engenharia Electrot\'ecnica e de Computadores, Faculdade de Ci\^encias e Tecnologia, FCT/UNL, Universidade Nova de Lisboa, 2829-516 Caparica, Portugal; and Instituto de Telecomunica\c{c}\~oes, 1049-001 Lisbon, Portugal (e-mail: pfa@fct.unl.pt)}%
  \\
  Diogo Raimundo%
  \thanks{Diogo Raimundo is with the Departamento de Engenharia Electrot\'ecnica e de Computadores, Faculdade de Ci\^encias e Tecnologia, FCT/UNL, Universidade Nova de Lisboa, 2829-516 Caparica, Portugal (e-mail: \{d.raimundo\}@campus.fct.unl.pt)}
}

\markboth{IEEE Transactions on Network and Service Management}%
{Amaral and Raimundo: Adversarial Robustness of MADDPG Routing}

\maketitle

\begin{abstract}
Learned routing policies expose a new attack surface: the observations the agents
read can be perturbed by a compromised telemetry channel. We study whether a
gradient-based observation attack (FGSM/PGD) can degrade the delivery performance of
Multi-Agent Deep Deterministic Policy Gradient (MADDPG) routing policies on a
realistic 86-node service-provider topology, across the same seven architectural
variants compared in a companion study. Our central finding separates two things
that adversarial-robustness studies routinely conflate: whether an attack changes a
policy's \emph{decisions}, and whether it changes the \emph{outcome}. Under an
$L_\infty$-bounded observation attack the FGSM gradient reliably flips routing
decisions --- up to $26\%$ of per-destination path choices at $\epsilon{=}0.3$,
roughly three times a budget-matched random perturbation --- yet delivery on the
intact, well-provisioned topology barely moves (at most $2.8$\,pp of an available
$21$\,pp damage ceiling). $K$-shortest-path redundancy converts changed decisions
into unchanged outcomes. Against a rigorous non-adversarial control (a budget- and
domain-matched \emph{random} perturbation) we find that the gradient's advantage is
real but small and architecture-dependent: four of seven variants show a
statistically significant gradient-over-random gap of $1$--$3$\,pp at nominal load,
while the simplest, best-delivering architecture is adversarially robust. Crucially,
under out-of-distribution link failures --- where redundancy is thin and delivery
does collapse under perturbation --- the random control is as damaging as, or more
damaging than, the gradient attack, so the failure-regime fragility is generic
observation-noise sensitivity rather than adversarial exploitability. We also find
that a graph-neural-network encoder suppresses decision flips for two of three base
architectures but not universally. The results argue that, for redundant transport
networks, telemetry integrity and fast reconvergence matter more than adversarial
training.
\end{abstract}
\begin{IEEEkeywords}
Adversarial machine learning; deep reinforcement learning; multi-agent systems;
adaptive routing; observation-space attacks; FGSM; robustness.
\end{IEEEkeywords}

\section{Introduction}
\label{sec:intro}
Deep reinforcement learning (DRL) is increasingly proposed for adaptive routing in
transport networks, and a companion study~\cite{companion} shows that a centralised-critic
MADDPG policy, trained at a concentrated-demand stress regime, outperforms a
shortest-path control plane and degrades gracefully under natural link failures.
That study concerns \emph{natural} stress. This paper asks the adjacent question:
does the same policy survive \emph{malicious} stress, specifically an adversary that
perturbs the observations the agents read?

Learned control planes consume telemetry --- per-path utilisation, queue occupancy,
link state --- to make decisions. If that telemetry channel is compromised, an
attacker can add a small, bounded perturbation to an agent's observation without
touching the network, the traffic, or the reward. This is the observation-space
threat model of adversarial machine learning~\cite{goodfellow2015,huang2017adversarial},
transplanted from image classification to a networked control loop. The natural
worry is that a policy which learned to exploit fine utilisation differences is
exactly a policy that a gradient attacker can fool.

Our contribution is to show that this worry, while mechanically valid, does not
translate into an operational vulnerability on a redundant transport topology --- and
to pin down precisely when it would. We make four points.
\begin{enumerate}
  \item \textbf{The attack works at the decision level.} FGSM/PGD flip a large
  fraction of routing decisions (up to $26\%$ at $\epsilon{=}0.3$), far more than a
  budget-matched random perturbation, confirming the gradient does targeted work
  (Section~\ref{sec:flips}).
  \item \textbf{The topology absorbs the flips.} Despite flipping a quarter of
  decisions, delivery falls by at most $2.8$\,pp against a damage ceiling of
  ${\sim}21$\,pp: $K$-path redundancy routes the moved flows onward
  (Section~\ref{sec:absorb}).
  \item \textbf{The adversarial signal is real but small and architecture-dependent.}
  Against a rigorous random control, four of seven variants show a significant
  $1$--$3$\,pp gradient-over-random gap; the simplest architecture and two of three
  GNN variants do not (Section~\ref{sec:arch}).
  \item \textbf{Failure-regime fragility is noise, not adversary.} When link
  failures remove redundancy and delivery does collapse under perturbation, a random
  perturbation is as damaging as the gradient --- so the collapse is observation-noise
  sensitivity, not adversarial exploitation (Section~\ref{sec:failures}).
\end{enumerate}

\section{Threat Model and Attack Methodology}
\label{sec:threat}

\subsection{System under attack}
We attack the seven MADDPG variants of the companion study without modification:
two critic domains (centralised, CC; local, LC) $\times$ two value heads (simple,
duelling) plus three graph-neural-network (GNN) encoder variants, all trained at the
$2\times$ concentrated-hotspot operating point on an 86-node service-provider
topology with a provider-edge, $K{=}3$ shortest-path forwarding model. Each agent,
for every destination, selects one of its $K$ candidate paths by an $\arg\max$ over
per-path logits. The decision is therefore \emph{discrete and per-destination}; the
attack must move a continuous observation far enough to cross an $\arg\max$ boundary.

The perturbed part of the observation is the per-path bottleneck (max-over-hops)
utilisation $U_i \in [0,1]^{|D|\times K}$ --- the exact signal that lets an agent
prefer a less-congested path, and therefore the natural target.

\subsection{Threat model}
The attacker adds a perturbation $\delta$ to an agent's observation, bounded in
$L_\infty$ norm by a budget $\epsilon$ (no single feature moves by more than
$\epsilon$), with bandwidth features re-projected into their valid $[0,1]$ range
after perturbation. The attacker cannot alter the network, the injected traffic, or
the reward signal. This models a compromised telemetry or monitoring channel and is
a white-box threat: the gradient attacks require access to the actor network.

\subsection{Attacks}
\textbf{FGSM.} We build a differentiable \emph{packet-loss objective} whose gradient
points toward worse routing: it weights each path by a sharpened function of its true
bottleneck utilisation and rewards placing policy probability mass on congested
paths, $\mathcal{L}(\delta) = \sum_{\text{paths}} \pi(\text{path})\,
\sigma\!\big((U-0.5)\cdot 10\big)$. The single-step Fast Gradient Sign Method
perturbs the observation by $\epsilon\cdot\operatorname{sign}(\nabla_{\text{obs}}
\mathcal{L})$, moving to the corner of the $\epsilon$-ball that most increases the
congestion objective. \textbf{PGD} iterates this step with re-projection into the
$\epsilon$-ball as a stronger upper bound; on this problem it flips a comparable
fraction of decisions to FGSM, because the $\arg\max$ margins to cross are small.

\textbf{Random control (the non-adversarial baseline).} The single most important
comparison in this paper is against a budget- and domain-matched \emph{random}
perturbation: same $\epsilon$, same clamping, but the direction is a random
$\pm\epsilon$ per feature rather than the gradient. This is the blind, zero-knowledge
member of the same attacker family; it is \emph{not} an adversarial attack. The
derived quantity we report throughout is the
\begin{equation}
\text{adversarial-specific effect} = \Delta_{\text{grad}} - \Delta_{\text{rand}},
\label{eq:gap}
\end{equation}
the difference between the delivery drop under the gradient attack and under the
random control. It isolates \emph{the value of gradient knowledge to the attacker}:
a value near zero means the attack is no better than blind noise; a positive value
is a genuine targeted vulnerability.

\subsection{Diagnostics and statistics}
Two diagnostics separate decisions from outcomes. The \textbf{action-flip rate} is
the fraction of per-(agent, destination) $\arg\max$ path choices that a perturbation
changes; it measures whether the attack moves the \emph{policy}. The \textbf{damage
ceiling} replaces the policy by fixed oracle routing rules on identical traffic:
the \emph{worst}-path rule (always the most-congested candidate) is the maximum
delivery loss any observation attacker could cause by steering decisions, bounding
what is at stake at each operating point. Evaluations use paired episodes: clean,
gradient, and random runs share the same seeded traffic and failure draws, so
per-episode differences are meaningful and the gap of Eq.~\eqref{eq:gap} carries a
proper paired $95\%$ confidence interval (CI); we use $15$ paired episodes per cell.

\section{Results}
\label{sec:results}

\subsection{FGSM changes routing decisions}
\label{sec:flips}
The attack is not under-powered. Figure~\ref{fig:flips-eps} reports the action-flip
rate against the perturbation budget for CC-Simple at the $2\times$ hotspot operating
point. The gradient flips an increasing share of decisions with $\epsilon$ ---
${\sim}8\%$ at $\epsilon{=}0.05$, ${\sim}14\%$ at $0.1$, ${\sim}22\%$ at $0.2$, and
$25.7\%$ at $0.3$ --- while a budget-matched random perturbation at $\epsilon{=}0.3$
flips only ${\sim}8\%$. The gradient thus flips roughly \emph{three times} as many
decisions as blind noise, confirming that it is doing targeted work: it locates and
crosses $\arg\max$ boundaries rather than jostling the observation. Across the
non-GNN variants the nominal flip rate is $18$--$26\%$. By the standard image-domain
reading, a policy whose decisions move this much is ``successfully attacked.''

\begin{figure}[t]
  \centering
  \includegraphics[width=0.92\linewidth]{T1_flip_vs_epsilon.pdf}
  \caption{Action-flip rate vs.\ perturbation budget $\epsilon$ (CC-Simple, $2\times$
  hotspot). FGSM (gradient) flips ${\sim}3\times$ more routing decisions than a
  budget-matched random perturbation, confirming the attack is targeted, not blind.}
  \label{fig:flips-eps}
\end{figure}

\subsection{Redundancy absorbs the flips: changed decisions, unchanged delivery}
\label{sec:absorb}
The outcome tells a different story from the decisions. Figure~\ref{fig:flips-pdr}
places, for every variant at nominal load, the fraction of decisions flipped beside
the delivery lost. Flipping $18$--$26\%$ of decisions costs \emph{at most} $2.8$\,pp
of end-to-end delivery (worst case LC-Simple), and $\approx 0$ for CC-Simple,
CC-Simple-GNN, and LC-Duelling. The mechanism is redundancy: a flipped decision moves
a flow from its chosen $K$-path to a different $K$-shortest path which, on a
well-provisioned topology with hop-diverse alternatives and spare capacity, still
delivers. The routing structure converts most changed decisions into identical
outcomes.

The damage ceiling quantifies how little of the available harm the attack extracts.
Figure~\ref{fig:ceiling} contrasts, across offered load, the clean policy delivery,
the \emph{worst}-path floor, and the greedy reference. At the $2\times$ operating
point roughly $21$\,pp of delivery separates the policy from the worst-path rule ---
that is what an omniscient decision-steering attacker could destroy --- yet FGSM
reaches at most $2.8$\,pp of it (${\leq}13\%$ of the ceiling), and $\approx 0$ on the
robust variants. The vulnerability exists in principle; the topology denies it in
practice.

\begin{figure}[t]
  \centering
  \includegraphics[width=0.92\linewidth]{T2_flips_vs_pdr_nominal.pdf}
  \caption{Decisions changed vs.\ delivery lost, all seven variants, $2\times$
  hotspot, $\epsilon{=}0.3$. Many decisions flip; almost no packets are lost.}
  \label{fig:flips-pdr}
\end{figure}

\begin{figure}[t]
  \centering
  \includegraphics[width=0.92\linewidth]{T7_damage_ceiling_contrast.pdf}
  \caption{Damage ceiling vs.\ achieved damage across offered load. The shaded band
  is the delivery an ideal decision-steering attacker could extract (policy minus
  worst-path); FGSM extracts almost none of it.}
  \label{fig:ceiling}
\end{figure}

\subsection{The adversarial signal is real, small, and architecture-dependent}
\label{sec:arch}
Absorbing the flips does not mean the gradient is worthless. Because clean, gradient,
and random runs are paired episode-for-episode, we can test the adversarial-specific
effect of Eq.~\eqref{eq:gap} with a proper CI. Figure~\ref{fig:gap} reports it for
all seven variants at nominal load. Four variants show a statistically significant
positive gap --- LC-Simple $+3.3$\,pp (CI $[+2.0,+4.6]$), CC-Duelling $+2.5$\,pp
($[+1.1,+3.9]$), LC-Duelling $+1.2$\,pp ($[+0.6,+1.9]$), and CC-Duelling-GNN
$+1.0$\,pp ($[0.0,+2.0]$) --- so on these architectures the gradient direction does
buy the attacker a real, if small, delivery loss over blind noise. The other three ---
CC-Simple ($+0.4$\,pp, CI $[-0.3,+1.0]$), CC-Simple-GNN ($\approx 0$), and
LC-Duelling-GNN ($\approx 0$) --- have gaps that include zero: on them the gradient
is no better than random. The pattern is consistent under a central critic the
\emph{duelling} head is more exploitable than the \emph{simple} head, and CC-Simple ---
the simplest and best-delivering architecture in the companion study --- is also among
the most adversarially robust. We stress the effect is minor operationally: absolute
loss is ${\leq}2.8$\,pp everywhere, so this is a ranking of a small effect, not a
vulnerability.

The GNN encoder mostly, but not universally, suppresses decision flips
(Figure~\ref{fig:gnn}). Two of the three GNN variants flip only $1.4\%$
(CC-Simple-GNN) and $2.6\%$ (LC-Duelling-GNN) of decisions --- roughly a tenth of their
non-GNN bases --- because message passing smooths the observation so that per-agent
perturbations rarely change an $\arg\max$. The exception is CC-Duelling-GNN, which
still flips $12.6\%$ and carries the small significant gap above: GNN decision
robustness is a property of the specific base architecture, not of GNN encoding as
such. Where it holds, the encoder that \emph{cost} delivery in the companion study
\emph{buys} decision robustness here --- a base-dependent security/performance
trade-off.

\begin{figure}[t]
  \centering
  \includegraphics[width=0.92\linewidth]{T3_adversarial_gap_by_variant.pdf}
  \caption{Adversarial-specific effect (gradient minus random delivery drop) at
  nominal load, seven variants, paired $95\%$ CI. Red bars exclude zero (a real
  targeted signal); grey bars include zero (no better than noise).}
  \label{fig:gap}
\end{figure}

\begin{figure}[t]
  \centering
  \includegraphics[width=0.92\linewidth]{T6_gnn_flip_robustness.pdf}
  \caption{Action-flip rate by variant (green: GNN). The GNN encoder suppresses flips
  for the Simple and LC-Duelling bases but not for CC-Duelling.}
  \label{fig:gnn}
\end{figure}

\subsection{Under failures, the fragility is noise, not adversary}
\label{sec:failures}
Redundancy is what absorbs the flips, so the sharpest test of adversarial
exploitability is to remove it. Figure~\ref{fig:fragility} sweeps the number of
random link failures per episode ($n=0,2,4,6$) at $2\times$ hotspot for CC-Simple,
plotting the delivery drop under the gradient attack and under the random control
with paired CI bands, and the clean (un-attacked) delivery for reference. Three
regimes appear. At $n{=}0$ both drops are negligible (Section~\ref{sec:absorb}). At
$n{=}2$ the network still delivers ${\sim}86\%$ cleanly, but now a perturbation
bites --- and it bites \emph{harder for random than for gradient}: the random control
drops delivery by ${\sim}6.5$\,pp while the gradient drops it by only ${\sim}0.3$\,pp,
a gap of $-6.1$\,pp whose CI excludes zero. At $n{=}4$ both attacks collapse delivery
together (${\sim}33$\,pp gradient, ${\sim}30$\,pp random); the $+3$\,pp gap is
\emph{not} significant, lost inside the failure-driven variance. At $n{=}6$ the clean
network has already self-collapsed to ${\sim}9\%$ delivery, so the attack-induced
drop shrinks simply because little delivery remains to lose --- the point is
failure-dominated and identical across all variants, and we exclude it from
attack conclusions.

The verdict is unambiguous: wherever the gradient's targeting could matter under
failure, a direction-free random perturbation is as damaging or more. The large
failure-regime delivery loss is therefore \emph{generic observation-noise
sensitivity}, not adversarial exploitability. Notably the flip rate \emph{rises}
under failure (Figure~\ref{fig:flip-fail}) --- post-failure states are more sensitive
and more decisions change --- yet, cross-referencing Figure~\ref{fig:fragility}, the
extra flips buy the adversary nothing over random noise: with few surviving paths any
perturbation that changes a decision can shove a flow onto an overloaded survivor,
and the gradient's structure confers no advantage.

\begin{figure}[t]
  \centering
  \includegraphics[width=0.92\linewidth]{T4_failure_fragility.pdf}
  \caption{Delivery drop vs.\ number of failed links (CC-Simple, $2\times$ hotspot),
  gradient and random with paired $95\%$ CI bands; grey dotted line is clean delivery
  (right axis). Fragility rises with failures, but random $\geq$ gradient throughout;
  the $n{=}6$ dip is network self-collapse, not renewed robustness.}
  \label{fig:fragility}
\end{figure}

\begin{figure}[t]
  \centering
  \includegraphics[width=0.92\linewidth]{T5_flip_vs_failures.pdf}
  \caption{Action-flip rate vs.\ number of failed links. The attack flips more
  decisions under failure, yet (cf.\ Fig.~\ref{fig:fragility}) gains no adversarial
  advantage over random noise.}
  \label{fig:flip-fail}
\end{figure}

\section{Discussion}
\label{sec:discussion}
The results are organised around a distinction that adversarial-robustness
evaluations often blur: an attack that changes a policy's \emph{decisions} is not the
same as an attack that changes the \emph{outcome}. On a redundant transport topology
these come apart cleanly. FGSM is a competent decision-level attack --- it flips a
quarter of routing choices, three times blind noise --- and yet the network's
$K$-path redundancy converts those changed decisions back into delivered packets.
The damage ceiling makes the point quantitatively: ${\sim}21$\,pp of delivery is at
stake, and the attack claims at most $2.8$\,pp of it.

The random control is what keeps the interpretation honest. Without it, the
$25$--$33$\,pp delivery collapse under a few link failures reads as a devastating
adversarial result; with it, the collapse is revealed as noise sensitivity, because a
direction-free perturbation of the same budget is equally or more damaging. We
recommend that observation-space robustness studies for networked control always
report a budget-matched random control alongside the gradient attack; the
gradient-minus-random gap, not the raw gradient drop, is the quantity that measures
adversarial vulnerability.

For operators the guidance is concrete. Adversarial training against FGSM would
defend a decision boundary that redundancy already protects, at nominal load, while
doing nothing for the failure regime, where the threat is undirected noise rather
than a clever adversary. Investment is better placed in \emph{telemetry integrity}
(so observations cannot be perturbed at all) and \emph{fast reconvergence} (so the
redundancy the policy relies on is restored quickly after failures). Architecturally,
the companion study's recommended CC-Simple policy is also the most adversarially
robust here, so the two studies agree on the same choice.

\section{Conclusion}
\label{sec:conclusion}
We asked whether a gradient observation attack can harm MADDPG routing, and found
that it changes many decisions but almost no outcomes on a redundant
service-provider topology: FGSM flips up to $26\%$ of routing decisions yet costs at
most $2.8$\,pp of a $21$\,pp damage ceiling. A rigorous random control shows the
gradient's advantage is real but small ($1$--$3$\,pp on four of seven architectures,
zero on the simplest and on two of three GNN variants), and that the large
failure-regime fragility is generic observation-noise sensitivity, not adversarial
exploitation. Limitations mirror the companion study: a single topology and one
family of attacks. Future work includes black-box and transfer attacks, attacks that
target the worst path directly to close the gap to the damage ceiling, and temporal
multi-step attacks that exploit the routing dynamics rather than a single
observation.

\begin{thebibliography}{1}
\bibitem{companion} P.~Amaral and D.~Raimundo, ``Centralised or Local Critic?
Mechanism-Grounded Selection of MADDPG Architectures for Adaptive Routing,''
companion manuscript, 2026.
\bibitem{goodfellow2015} I.~J. Goodfellow, J.~Shlens, and C.~Szegedy, ``Explaining
and Harnessing Adversarial Examples,'' in \emph{Proc. ICLR}, 2015.
\bibitem{huang2017adversarial} S.~Huang, N.~Papernot, I.~Goodfellow, Y.~Duan, and
P.~Abbeel, ``Adversarial Attacks on Neural Network Policies,'' in \emph{Proc. ICLR
Workshop}, 2017.
\end{thebibliography}

\end{document}
